\section{Loi aréale de la coupe triadique : capacité \(81\log 3\), opérateur d’aire, spectre et entrelaceur collectif des canaux \(90/120\)}
\label{hap:sec:area-entrelazamiento-barbero-rev11}
\index{Barbero--Immirzi!opérateur d’aire}
\index{entropie modulaire!Hamiltonien modulaire}

L’identification HMT de la fonctionnelle \(\Gamma_{6\to8}\) à la coordonnée
de Barbero--Immirzi s’effectue sur le même bord triadique qui porte
l’incidence (90/120). L’intérieur et son feuillet de bord ont la même cardinalité,
\[
 \mathcal B=(\mathbb Z/9\mathbb Z)^3,
 \quad |\mathcal B|=729,
 \qquad
 \partial\mathcal B=(\mathbb Z/27\mathbb Z)^2,
 \quad |\partial\mathcal B|=729,
\]
par le regroupement de six trits en trois coordonnées nonadiques ou deux
coordonnées d’ordre vingt-sept.

\begin{theorem}[Loi aréale de la coupe triadique]
\label{hap:thm:ley-areal-corte-triadico-rev11}
Dans le graphe cubique \(N\times N\times N\) de capacité unitaire, la coupe
minimale entre deux faces opposées a pour capacité \(N^2\). Pour \(N=9\),
\[
 \mathcal A_{\min}=81,
 \qquad S_{\rm tri}=\mathcal A_{\min}\log3=81\log3.
\]
\end{theorem}
\begin{proof}
Les \(N^2\) droites parallèles à la direction normale sont des chemins disjoints en
arêtes, de sorte que toute coupe les intercepte. Une couche transversale contient
exactement \(N^2\) arêtes et atteint la borne. Chaque liaison triadique apporte
\(\log3\) à l’entropie d’un état maximalement mélangé.
\end{proof}

Dans le tore de bord \(27\times27\), soit \(A\) le bloc canonique
\(9\times9\). Sa frontière se décompose exactement en
\begin{equation}
 \partial A=\partial A_{90}\sqcup\partial A_{120},
 \qquad |\partial A_{90}|=|\partial A_{120}|=18.
 \label{hap:eq:corte-90-120-rev11}
\end{equation}
Les liaisons horizontales représentent le canal (90), et les verticales le
canal (120). Cette décomposition transforme la paire d’incidences en une
décomposition opératorielle du bord.

L’entropie aréale \(S_{\rm tri}=81\log3\) appartient à l’état maximalement
mélangé des liaisons de la coupe. L’entropie bicouche
\(S_{\rm bi}(y)\) de la \cref{hap:def:entropia-bicapa-rev6} appartient à la
distribution normalisée entre les deux feuillets ; les deux objets conservent
des domaines distincts.

Soit \(N=\operatorname{diag}(0,1,2)\) l’opérateur nombre d’une liaison tritique.
La réalisation aréale déclare l’échelle positive
\[
 \theta_{\rm clk}=\frac{C^*}{6}\frac{\pi}{180},
 \qquad
 a_0=\frac{1+2\cos(10^\circ)}{3}\,\theta_{\rm clk},
 \qquad \gamma_{\rm HMT}=\Gamma_{6\to8}.
\]
La formule de \(a_0\) constitue la structure de départ de cette
réalisation ; une fois cette échelle fixée, l’opérateur et son spectre se déduisent
exactement.

\begin{definition}[Opérateur d’aire HMT]
Sur le feuillet de bord
\(\mathcal H_{\partial}=\bigotimes_{e\in\partial A}\mathbb C^3\), on définit
\begin{equation}
 \widehat{\mathcal A}_{\rm HMT}
 =\gamma_{\rm HMT}a_0\sum_{e\in\partial A}N_e.
 \label{hap:eq:operador-area-hmt-rev11}
\end{equation}
Son spectre dans la coupe canonique est
\[
 \operatorname{Spec}\widehat{\mathcal A}_{\rm HMT}
 =\{n\gamma_{\rm HMT}a_0:0\leq n\leq72\},
\]
avec des multiplicités données par les coefficients de
\((1+z+z^2)^{36}\). Le prolongement par coupes compatibles produit la
famille protégée \(\mathcal A_n^{\rm prot}=n\gamma_{\rm HMT}a_0\).
\end{definition}

L’entrelaceur qui fixe la normalisation se construit dans le sous-espace des
canaux collectifs. Soient \(u_{90},u_{120}\) les vecteurs unitaires uniformes
des incidences hexadique et octadique, et soient \(v_{90},v_{120}\) les
vecteurs unitaires uniformes des deux ensembles de liaisons de
\eqref{hap:eq:corte-90-120-rev11}. L’application
\begin{equation}
 \begin{aligned}
 \mathcal J_{6\to8}:{}
 &\operatorname{span}\{u_{90},u_{120}\}
 \longrightarrow
 \operatorname{span}\{v_{90},v_{120}\},\\
 &\mathcal J_{6\to8}u_{90}=v_{90},
 \qquad
 \mathcal J_{6\to8}u_{120}=v_{120}
 \end{aligned}
 \label{hap:eq:entrelazador-area-barbero-rev11}
\end{equation}
est une isométrie entre ces sous-espaces collectifs. Si
\[
 D_{\rm inc}=90|u_{90}\rangle\langle u_{90}|
 +120|u_{120}\rangle\langle u_{120}|,
 \qquad
 D_{\rm area}=90|v_{90}\rangle\langle v_{90}|
 +120|v_{120}\rangle\langle v_{120}|,
\]
alors
\[
 \boxed{\mathcal J_{6\to8}D_{\rm inc}
 =D_{\rm area}\mathcal J_{6\to8}.}
\]
En particulier, la combinaison primitive \(12:-1\) agit
avec les mêmes coefficients avant et après \(\mathcal J_{6\to8}\). Cette
égalité fixe le dictionnaire de normalisation qui insère
\(\gamma_{\rm HMT}=\Gamma_{6\to8}\) dans la connexion
d’Ashtekar--Barbero.

